\documentclass{article}
\usepackage[utf8]{inputenc}
\usepackage{geometry}
\usepackage{enumitem}
\usepackage{hyperref}
\usepackage{xcolor}
\usepackage{graphicx}
\usepackage{array}
\usepackage{multirow}
\usepackage{amsmath, amssymb}
\usepackage{chemfig}
\usepackage{mhchem}
\usepackage{tikz}
\usepackage{setspace}
\usepackage{soul}
\usepackage{mathtools}
\usepackage{booktabs}
\usepackage{chemformula}
\usepackage{siunitx}
\usetikzlibrary{positioning, calc}
\usepackage{longtable}
\geometry{margin=1in}
\setlength{\parskip}{0.5em}
\usepackage{cancel}

\begin{document}

\begin{center}
    {\LARGE\textbf{SI Session Planning Form}}
\end{center}

\
\noindent
\begin{flushleft}
\textbf{SI Leader:} Brandon Cobb\\
\textbf{Session Week:} 7\\
\textbf{Session\#:} 7\\
\textbf{Instructor:} Professor Tramontozzi\\
\textbf{Old Plans:}\href{https://macombedu.sharepoint.com/:b:/s/ASC-SupplementalInstructionEmbeddedTutoring/EeSYqCwT_KhMolgJvXbVEjgBOyF_DevvS8BgCm1zPkQQJg?e=QJTt4o}{SharePoint}\\
\textbf{Session Goal:} Students will have reviewed balancing chemical equations and stoichiometry, applying past concepts like formulas, atomic masses, avogadro's number etc. to determine a quanitity associated with a given chemical or molar equivalent to a different chemical, likely in a reaction. 
\end{flushleft}

\vspace{1em}
\section*{\underline{Opener}}
{\centering\large\textbf{\hl{Take Attendance}}\par}\vspace{-0.25\baselineskip}
\noindent
\begin{tabular}{|p{4.2cm}|p{9.8cm}|}
\hline
\textbf{Collaborative Learning Techniques} & Group Discussion\\
\hline
\textbf{Learning Strategy} & Lecture Review\\
\hline
\textbf{Time} & 12:00---12:10\\
\hline
\textbf{Materials} & Notes, homework, whiteboard, markers, ready to speak, periodic table \\
\hline
\textbf{Lecture/Textbook/Old Plans/Slide References} & Class roster, Chapter 6: Chemical calculations: Formula masses, moles and chemical equations.  \\
\hline
\end{tabular}


\vspace{0.5cm}
{\large\textbf{Definitions}}
\\
\begin{tabular}{|p{0.9\linewidth}|}
\hline
\begin{minipage}[t]{0.9\linewidth}
\begin{enumerate}
    \item\textbf{Group Discussion:} One person in the group is asked to present on a topic or review material for the group and then lead the discussion for the group. This person should not be the regular group leader.
   \item\textbf{Lecture Review:} As a group summarize the lecture from the previous class. You may have to provide prompts for the students. For example, "The first concept discussed was Civil Liberties adn Public Policy, what did the profressor highlight regarding this?" You may want to ask them to try summarizing without looking at their notes; however, if they are having a difficult time remembering, tell them to refer to their notes.
\end{enumerate}
\end{minipage} \\
\hline
\end{tabular}
\newpage
\noindent
{\large\textbf{Instructions (please provide a\hl{DETAILED} activity breakdown below (and/or attached}}\\

\begin{tabular}{|p{0.9\linewidth}|}
\hline
\vspace{0.2cm}
\textbf{Lecture Review:}
\begin{enumerate}
\item During te first 10-15 minutes of the SI session. have the students summarize the most recent lecture, or have them identify the key words from that lecture.
\item Give students three minutes to find support in their lecture notes for a given generalization.
\item Have the students predict the direction of future lectures based upon the past lectures.
\item Have the student arrange terms fgrom lecture and text into a structured outline.
\item Reinforce new terms or important information by using clearly constructed handouts (can be complete or nearly complete at the beginning of the term but should gradually require more and more filling in as the group becomes more accustomed to working together.
\item Review material from previous sessions and lectures.
\item Take a couple of minutes at the end of the SI session to summarize the main idea convered during a session. Asl the students to help summarize.
\item Have students write a one paragraph summary of the lecture. List the new vocabulary terms introduced with this lecture.
\item Formulate potential exam questions based on the main ideas from the lecture.
\item Formulate pontential answers from details in the lecture notes.
\end{enumerate} \\
\hline
\end{tabular}

\newpage
\noindent
{\large\textbf{Checking for Understanding (questions \& answers) continued...} \\
\vspace{0.2cm}
\begin{tabular}{|p{0.9\linewidth}|}
\hline
\begin{enumerate}
\item Why do chemical equations need to be balanced in the first place?
\item What does it mean when an equation is balanced?
\item Why do we use coefficients to balance equations instead of changing subscripts?
\item What does a balanced equation tell us about the number of particles reacting?
\item How does a balanced equation connect to the law of conservation of mass?
\item Why can we use moles to compare amounts of different substances in a reaction?
\item Why do chemists use the mole instead of counting individual atoms or molecules?
\item What does “one mole” actually represent?
\item Why does one mole of carbon weigh differently than one mole of oxygen?
\item How is molar mass related to the periodic table?
\item Why do we say molar mass acts as a “bridge” between grams and moles?
\item Why can we use molar mass to find the number of particles in a sample?
\item Why does the formula mass of a compound depend on how many atoms of each element it contains?
\item Why do compounds with different formulas have different molar masses?
\item What does the formula of a compound tell us about the ratio of atoms inside it?
\item Why can we predict how much product will form just by knowing the balanced equation and moles?
\item Why does the ratio of coefficients in a balanced equation matter for calculations?
\item Why is understanding units important when solving mole problems?
\item How does knowing the molar mass help convert between substances in real experiments?
\item Why do chemists always check that equations are balanced before doing stoichiometry?
\end{enumerate} \\
\hline
\end{tabular}


\section*{\underline{Activity 1}}
\vspace{0.2cm}
\begin{tabular}{|p{4.2cm}|p{9.8cm}|}
\hline
\textbf{Collaborative Learning Techniques} & Jigsaw\\
\hline
\textbf{Learning Strategy} & Predict Test Questions \\
\hline
\textbf{Time} & 12:10---12:30\\
\hline
\textbf{Materials} & Pencil, eraser, calculator, periodic table \\
\hline
\textbf{Lecture/Textbook/Old Plans/Slide References} & Chapter 6: Chemical calculations: Formula masses, moles and chemical equations \\
\hline
\end{tabular}
\vspace{0.5cm}

{\large\textbf{Definition}}\\
\noindent
\begin{tabular}{|p{0.9\linewidth}|}
\hline
\begin{enumerate}
    \item\textbf{Jigsaw:} Jigsaws, when used properly, make the group as a whole dependent upon all of them in subgroups. Each group provides a peice of the puzzle. Group members are broken into smaller groups. Each small group works on some aspect of the same problem, question, or issue. They then share their part of the puzzle with the large group.
    \item\textbf{Predict Test Questions:}  Put students in groups of 2 or 3 and assign them to write a test question for a specific topic, ensuring that all topics have been convered. Ask students to write their questions on the board or on an overhead for discussion (would the professor ask the question? What is the answer? etc.). Students will have the benefit of learning to think like the teacher and they'll be able to see additional questions that other students have writen.
 
\end{enumerate}\\
\hline
\end{tabular}
\vspace{0.2cm}
\newpage
{\large\textbf{Instructions (please provide a\hl{DETAILED} activity breakdown below (and/or attached}}\\
\noindent
\begin{tabular}{|p{0.9\linewidth}|}
\hline
\vspace{0.2cm}
\begin{enumerate}
    \item \textbf{Distribute Worksheets:} Each student receives a worksheet containing chemistry questions. The problems cover a range of difficulty so that students can decide which ones best match their needs.
    
    \item \textbf{Student Choice of Starting Point:} Students look over the worksheet and select their own starting problem rather than being assigned a fixed order. This allows them to challenge themselves where they feel it is most beneficial.
    
    \item \textbf{Staggered Timing:} Students work in pairs. Each partner begins at a different point (for example, one may start at Problem 1 while the other starts at Problem 3). They continue by alternating every three questions. This staggered structure ensures that students will be asking for help or requesting solutions at different times, keeping the class flow balanced.
    
    \item \textbf{Independent Work:} Students attempt their chosen problems individually, writing down both reasoning and answers. Emphasis is placed on documenting thought process as well as the result.
    
    \item \textbf{Help and Solution Reveal:} After the set time for each problem expires, the SI leader either walks through the solution on the whiteboard or directs students to check the answer key. Students are encouraged to ask clarifying questions at the timed checkpoints rather than all at once.
    
    \item \textbf{Partner Collaboration:} Once a solution is reviewed, partners share their reasoning for their respective problems. They ask each other questions, offer corrections, and compare how their approaches aligned with the solution.
    
    \item \textbf{Rotation of Problems:} After every three questions, students switch to new problems, continuing the staggered order so that they stay offset from each other. This ensures that each student engages with a variety of problems across the worksheet.
    
    \item \textbf{Class-Wide Debrief:} The SI leader pauses periodically to highlight frequently asked questions, difficult concepts, or common errors observed during the staggered work. Whole-class solutions are discussed when many students identify the same challenge.
    
    \item \textbf{Reflection:} At the close, students note which problems gave them the most insight, what strategies helped, and what topics they still want clarified. This reflection helps focus future review sessions.
\end{enumerate}\\
\hline
\end{tabular}
\newpage
\noindent
{\large\textbf{Content (fully worked out problems \& answers) continued...}}\\[0.2em]
\vspace{0.2cm}
\noindent \begin{tabular}{|p{0.9\linewidth}|}
\hline
\vspace{0.2cm}
\textbf{(1)} Balance: \(\mathrm{H_2 + O_2 \rightarrow H_2O}\)\\
\begin{minipage}[t]{0.85\linewidth}
\begin{enumerate}[label=\alph*.]
   \item \(\mathrm{H_2 + O_2 \rightarrow 2H_2O}\)
   \item \(\mathrm{4H_2 + 2O_2 \rightarrow 4H_2O}\)
   \item \(\mathrm{2H_2 + O_2 \rightarrow 2H_2O}\)
   \item \(\mathrm{H_2 + \tfrac{1}{2}O_2 \rightarrow H_2O}\)
\end{enumerate}
\end{minipage}\\
\hline
\end{tabular}

\vspace{0.2cm}
\noindent \begin{tabular}{|p{0.9\linewidth}|}
\hline
\vspace{0.2cm}
\textbf{(2)} Balance: \(\mathrm{C_3H_8 + O_2 \rightarrow CO_2 + H_2O}\)\\
\begin{minipage}[t]{0.85\linewidth}
\begin{enumerate}[label=\alph*.]
   \item \(\mathrm{2C_3H_8 + 7O_2 \rightarrow 6CO_2 + 8H_2O}\)
   \item \(\mathrm{C_3H_8 + 3O_2 \rightarrow 3CO_2 + 4H_2O}\)
   \item \(\mathrm{C_3H_8 + 5O_2 \rightarrow 3CO_2 + 4H_2O}\)
   \item \(\mathrm{C_3H_8 + 4O_2 \rightarrow 3CO_2 + 4H_2O}\)
\end{enumerate}
\end{minipage}\\
\hline
\end{tabular}

\vspace{0.2cm}
\noindent \begin{tabular}{|p{0.9\linewidth}|}
\hline
\vspace{0.2cm}
\textbf{(3)} Balance: \(\mathrm{Fe + O_2 \rightarrow Fe_2O_3}\)\\
\begin{minipage}[t]{0.85\linewidth}
\begin{enumerate}[label=\alph*.]
   \item \(\mathrm{2Fe + O_2 \rightarrow Fe_2O_3}\)
   \item \(\mathrm{Fe + \tfrac{3}{2}O_2 \rightarrow Fe_2O_3}\)
   \item \(\mathrm{4Fe + 3O_2 \rightarrow 2Fe_2O_3}\)
   \item \(\mathrm{4Fe + 6O_2 \rightarrow 4Fe_2O_3}\)
\end{enumerate}
\end{minipage}\\
\hline
\end{tabular}

\vspace{0.2cm}
\noindent \begin{tabular}{|p{0.9\linewidth}|}
\hline
\vspace{0.2cm}
\textbf{(4)} Balance: \(\mathrm{N_2 + H_2 \rightarrow NH_3}\)\\
\begin{minipage}[t]{0.85\linewidth}
\begin{enumerate}[label=\alph*.]
   \item \(\mathrm{N_2 + H_2 \rightarrow 2NH_3}\)
   \item \(\mathrm{N_2 + \tfrac{3}{2}H_2 \rightarrow NH_3}\)
   \item \(\mathrm{2N_2 + 3H_2 \rightarrow 4NH_3}\)
   \item \(\mathrm{N_2 + 3H_2 \rightarrow 2NH_3}\)
\end{enumerate}
\end{minipage}\\
\hline
\end{tabular}
\newpage

\vspace{0.2cm}
\noindent \begin{tabular}{|p{0.9\linewidth}|}
\hline
\vspace{0.2cm}
\textbf{(5)} Balance: \(\mathrm{KClO_3 \rightarrow KCl + O_2}\) (thermal decomposition)\\
\begin{minipage}[t]{0.85\linewidth}
\begin{enumerate}[label=\alph*.]
   \item \(\mathrm{KClO_3 \rightarrow KCl + O_2}\)
   \item \(\mathrm{2KClO_3 \rightarrow KCl + 3O_2}\)
   \item \(\mathrm{2KClO_3 \rightarrow 2KCl + 3O_2}\)
   \item \(\mathrm{4KClO_3 \rightarrow 4KCl + 3O_2}\)
\end{enumerate}
\end{minipage}\\
\hline
\end{tabular}

\vspace{0.2cm}
\noindent \begin{tabular}{|p{0.9\linewidth}|}
\hline
\vspace{0.2cm}
\textbf{(6)} Balance: \(\mathrm{Al + O_2 \rightarrow Al_2O_3}\)\\
\begin{minipage}[t]{0.85\linewidth}
\begin{enumerate}[label=\alph*.]
   \item \(\mathrm{2Al + O_2 \rightarrow Al_2O_3}\)
   \item \(\mathrm{4Al + 3O_2 \rightarrow 2Al_2O_3}\)
   \item \(\mathrm{Al + \tfrac{3}{2}O_2 \rightarrow Al_2O_3}\)
   \item \(\mathrm{4Al + 6O_2 \rightarrow 4Al_2O_3}\)
\end{enumerate}
\end{minipage}\\
\hline
\end{tabular}

\vspace{0.2cm}
\noindent \begin{tabular}{|p{0.9\linewidth}|}
\hline
\vspace{0.2cm}
\textbf{(7)} Balance: \(\mathrm{CuO + H_2 \rightarrow Cu + H_2O}\)\\
\begin{minipage}[t]{0.85\linewidth}
\begin{enumerate}[label=\alph*.]
   \item \(\mathrm{CuO + 2H_2 \rightarrow Cu + 2H_2O}\)
   \item \(\mathrm{CuO + H_2 \rightarrow Cu + H_2O}\)
   \item \(\mathrm{2CuO + 2H_2 \rightarrow 2Cu + 2H_2O}\)
   \item \(\mathrm{2CuO + H_2 \rightarrow 2Cu + H_2O}\)
\end{enumerate}
\end{minipage}\\
\hline
\end{tabular}

\vspace{0.2cm}
\noindent \begin{tabular}{|p{0.9\linewidth}|}
\hline
\vspace{0.2cm}
\textbf{(8)} Balance: \(\mathrm{Zn + HCl \rightarrow ZnCl_2 + H_2}\)\\
\begin{minipage}[t]{0.85\linewidth}
\begin{enumerate}[label=\alph*.]
   \item \(\mathrm{2Zn + HCl \rightarrow ZnCl_2 + H_2}\)
   \item \(\mathrm{Zn + HCl \rightarrow ZnCl_2 + 2H_2}\)
   \item \(\mathrm{Zn + 2HCl \rightarrow ZnCl_2 + H_2}\)
   \item \(\mathrm{Zn + 2HCl \rightarrow ZnCl + H_2}\)
\end{enumerate}
\end{minipage}\\
\hline
\end{tabular}
\newpage

\vspace{0.2cm}
\noindent \begin{tabular}{|p{0.9\linewidth}|}
\hline
\vspace{0.2cm}
\textbf{(9)} Balance: \(\mathrm{Na + H_2O \rightarrow NaOH + H_2}\)\\
\begin{minipage}[t]{0.85\linewidth}
\begin{enumerate}[label=\alph*.]
   \item \(\mathrm{Na + 2H_2O \rightarrow NaOH + H_2}\)
   \item \(\mathrm{2Na + 2H_2O \rightarrow 2NaOH + H_2}\)
   \item \(\mathrm{Na + H_2O \rightarrow NaOH + H_2}\)
   \item \(\mathrm{2Na + H_2O \rightarrow 2NaOH + H_2}\)
\end{enumerate}
\end{minipage}\\
\hline
\end{tabular}

\vspace{0.2cm}
\noindent \begin{tabular}{|p{0.9\linewidth}|}
\hline
\vspace{0.2cm}
\textbf{(10)} Balance: \(\mathrm{AgNO_3 + NaCl \rightarrow AgCl + NaNO_3}\)\\
\begin{minipage}[t]{0.85\linewidth}
\begin{enumerate}[label=\alph*.]
   \item \(\mathrm{2AgNO_3 + NaCl \rightarrow AgCl + 2NaNO_3}\)
   \item \(\mathrm{2AgNO_3 + 2NaCl \rightarrow 2AgCl + 2NaNO_3}\)
   \item \(\mathrm{AgNO_3 + NaCl \rightarrow AgCl + NaNO_3}\)
   \item \(\mathrm{AgNO_3 + 2NaCl \rightarrow AgCl + 2NaNO_3}\)
\end{enumerate}
\end{minipage}\\
\hline
\end{tabular}

\vspace{0.2cm}
\noindent \begin{tabular}{|p{0.9\linewidth}|}
\hline
\vspace{0.2cm}
\textbf{(11)} Balance: \(\mathrm{HCl + NaOH \rightarrow NaCl + H_2O}\)\\
\begin{minipage}[t]{0.85\linewidth}
\begin{enumerate}[label=\alph*.]
   \item \(\mathrm{2HCl + NaOH \rightarrow NaCl + 2H_2O}\)
   \item \(\mathrm{HCl + NaOH \rightarrow NaCl + H_2O}\)
   \item \(\mathrm{HCl + 2NaOH \rightarrow NaCl + 2H_2O}\)
   \item \(\mathrm{2HCl + 2NaOH \rightarrow 2NaCl + 2H_2O}\)
\end{enumerate}
\end{minipage}\\
\hline
\end{tabular}

\vspace{0.2cm}
\noindent \begin{tabular}{|p{0.9\linewidth}|}
\hline
\vspace{0.2cm}
\textbf{(12)} Balance: \(\mathrm{C_6H_{12}O_6 + O_2 \rightarrow CO_2 + H_2O}\) (combustion)\\
\begin{minipage}[t]{0.85\linewidth}
\begin{enumerate}[label=\alph*.]
   \item \(\mathrm{C_6H_{12}O_6 + 3O_2 \rightarrow 6CO_2 + 6H_2O}\)
   \item \(\mathrm{2C_6H_{12}O_6 + 12O_2 \rightarrow 12CO_2 + 12H_2O}\)
   \item \(\mathrm{C_6H_{12}O_6 + 9O_2 \rightarrow 6CO_2 + 6H_2O}\)
   \item \(\mathrm{C_6H_{12}O_6 + 6O_2 \rightarrow 6CO_2 + 6H_2O}\)
\end{enumerate}
\end{minipage}\\
\hline
\end{tabular}
\newpage

\vspace{0.2cm}
\noindent \begin{tabular}{|p{0.9\linewidth}|}
\hline
\vspace{0.2cm}
\textbf{(13)} Balance: \(\mathrm{CaO + CO_2 \rightarrow CaCO_3}\)\\
\begin{minipage}[t]{0.85\linewidth}
\begin{enumerate}[label=\alph*.]
   \item \(\mathrm{CaO + 2CO_2 \rightarrow CaCO_3}\)
   \item \(\mathrm{CaO + CO_2 \rightarrow CaCO_3}\)
   \item \(\mathrm{2CaO + 2CO_2 \rightarrow 2CaCO_3}\)
   \item \(\mathrm{CaO + CO_2 \rightarrow 2CaCO_3}\)
\end{enumerate}
\end{minipage}\\
\hline
\end{tabular}

\vspace{0.2cm}
\noindent \begin{tabular}{|p{0.9\linewidth}|}
\hline
\vspace{0.2cm}
\textbf{(14)} Balance: \(\mathrm{Mg + N_2 \rightarrow Mg_3N_2}\)\\
\begin{minipage}[t]{0.85\linewidth}
\begin{enumerate}[label=\alph*.]
   \item \(\mathrm{2Mg + N_2 \rightarrow Mg_3N_2}\)
   \item \(\mathrm{6Mg + 2N_2 \rightarrow 2Mg_3N_2}\)
   \item \(\mathrm{3Mg + N_2 \rightarrow Mg_3N_2}\)
   \item \(\mathrm{Mg + N_2 \rightarrow Mg_3N_2}\)
\end{enumerate}
\end{minipage}\\
\hline
\end{tabular}

\vspace{0.2cm}
\noindent \begin{tabular}{|p{0.9\linewidth}|}
\hline
\vspace{0.2cm}
\textbf{(15)} Balance: \(\mathrm{Na_2CO_3 + HCl \rightarrow NaCl + H_2O + CO_2}\)\\
\begin{minipage}[t]{0.85\linewidth}
\begin{enumerate}[label=\alph*.]
   \item \(\mathrm{Na_2CO_3 + 3HCl \rightarrow 2NaCl + H_2O + CO_2}\)
   \item \(\mathrm{Na_2CO_3 + HCl \rightarrow NaCl + H_2O + CO_2}\)
   \item \(\mathrm{Na_2CO_3 + 2HCl \rightarrow 2NaCl + H_2O + CO_2}\)
   \item \(\mathrm{2Na_2CO_3 + 4HCl \rightarrow 4NaCl + 2H_2O + 2CO_2}\)
\end{enumerate}
\end{minipage}\\
\hline
\end{tabular}

\vspace{0.2cm}
\noindent \begin{tabular}{|p{0.9\linewidth}|}
\hline
\vspace{0.2cm}
\textbf{(16)} Balance: \(\mathrm{H_2O_2 \rightarrow H_2O + O_2}\) (decomposition)\\
\begin{minipage}[t]{0.85\linewidth}
\begin{enumerate}[label=\alph*.]
   \item \(\mathrm{H_2O_2 \rightarrow 2H_2O + O_2}\)
   \item \(\mathrm{4H_2O_2 \rightarrow 4H_2O + 2O_2}\)
   \item \(\mathrm{2H_2O_2 \rightarrow 2H_2O + O_2}\)
   \item \(\mathrm{H_2O_2 \rightarrow H_2O + O_2}\)
\end{enumerate}
\end{minipage}\\
\hline
\end{tabular}
\newpage

\vspace{0.2cm}
\noindent \begin{tabular}{|p{0.9\linewidth}|}
\hline
\vspace{0.2cm}
\textbf{(17)} Balance: \(\mathrm{C_2H_5OH + O_2 \rightarrow CO_2 + H_2O}\) (ethanol combustion)\\
\begin{minipage}[t]{0.85\linewidth}
\begin{enumerate}[label=\alph*.]
   \item \(\mathrm{C_2H_5OH + 6O_2 \rightarrow 2CO_2 + 3H_2O}\)
   \item \(\mathrm{2C_2H_5OH + 6O_2 \rightarrow 4CO_2 + 6H_2O}\)
   \item \(\mathrm{C_2H_5OH + \tfrac{3}{2}O_2 \rightarrow 2CO_2 + 3H_2O}\)
   \item \(\mathrm{C_2H_5OH + 3O_2 \rightarrow 2CO_2 + 3H_2O}\)
\end{enumerate}
\end{minipage}\\
\hline
\end{tabular}

\vspace{0.2cm}
\noindent \begin{tabular}{|p{0.9\linewidth}|}
\hline
\vspace{0.2cm}
\textbf{(18)} Balance: \(\mathrm{SO_2 + O_2 \rightarrow SO_3}\)\\
\begin{minipage}[t]{0.85\linewidth}
\begin{enumerate}[label=\alph*.]
   \item \(\mathrm{SO_2 + \tfrac{1}{2}O_2 \rightarrow SO_3}\)
   \item \(\mathrm{4SO_2 + 2O_2 \rightarrow 4SO_3}\)
   \item \(\mathrm{2SO_2 + O_2 \rightarrow 2SO_3}\)
   \item \(\mathrm{SO_2 + O_2 \rightarrow SO_3}\)
\end{enumerate}
\end{minipage}\\
\hline
\end{tabular}

\vspace{0.2cm}
\noindent \begin{tabular}{|p{0.9\linewidth}|}
\hline
\vspace{0.2cm}
\textbf{(19)} Balance: \(\mathrm{SO_3 + H_2O \rightarrow H_2SO_4}\)\\
\begin{minipage}[t]{0.85\linewidth}
\begin{enumerate}[label=\alph*.]
   \item \(\mathrm{2SO_3 + H_2O \rightarrow H_2SO_4}\)
   \item \(\mathrm{SO_3 + 2H_2O \rightarrow H_2SO_4}\)
   \item \(\mathrm{2SO_3 + 2H_2O \rightarrow 2H_2SO_4}\)
   \item \(\mathrm{SO_3 + H_2O \rightarrow H_2SO_4}\)
\end{enumerate}
\end{minipage}\\
\hline
\end{tabular}

\vspace{0.2cm}
\noindent \begin{tabular}{|p{0.9\linewidth}|}
\hline
\vspace{0.2cm}
\textbf{(20)} Balance: \(\mathrm{CH_4 + Cl_2 \rightarrow CH_3Cl + HCl}\) (chlorination, first step)\\
\begin{minipage}[t]{0.85\linewidth}
\begin{enumerate}[label=\alph*.]
   \item \(\mathrm{2CH_4 + Cl_2 \rightarrow CH_3Cl + HCl}\)
   \item \(\mathrm{CH_4 + 2Cl_2 \rightarrow CH_3Cl + 2HCl}\)
   \item \(\mathrm{CH_4 + Cl_2 \rightarrow CH_3Cl + HCl}\)
   \item \(\mathrm{2CH_4 + 2Cl_2 \rightarrow 2CH_3Cl + 2HCl}\)
\end{enumerate}
\end{minipage}\\
\hline
\end{tabular}
\newpage

\vspace{0.2cm}
\noindent \begin{tabular}{|p{0.9\linewidth}|}
\hline
\vspace{0.2cm}
\textbf{(21)} Balance: \(\mathrm{C_2H_4 + H_2 \rightarrow C_2H_6}\) (hydrogenation)\\
\begin{minipage}[t]{0.85\linewidth}
\begin{enumerate}[label=\alph*.]
   \item \(\mathrm{C_2H_4 + \tfrac{1}{2}H_2 \rightarrow C_2H_6}\)
   \item \(\mathrm{C_2H_4 + 2H_2 \rightarrow C_2H_6}\)
   \item \(\mathrm{2C_2H_4 + H_2 \rightarrow 2C_2H_6}\)
   \item \(\mathrm{C_2H_4 + H_2 \rightarrow C_2H_6}\)
\end{enumerate}
\end{minipage}\\
\hline
\end{tabular}

\vspace{0.2cm}
\noindent \begin{tabular}{|p{0.9\linewidth}|}
\hline
\vspace{0.2cm}
\textbf{(22)} Balance: \(\mathrm{NH_4NO_3 \rightarrow N_2O + H_2O}\) (thermal decomposition)\\
\begin{minipage}[t]{0.85\linewidth}
\begin{enumerate}[label=\alph*.]
   \item \(\mathrm{NH_4NO_3 \rightarrow N_2 + 2H_2O}\)
   \item \(\mathrm{2NH_4NO_3 \rightarrow 2N_2O + 2H_2O}\)
   \item \(\mathrm{NH_4NO_3 \rightarrow N_2O + H_2O}\)
   \item \(\mathrm{NH_4NO_3 \rightarrow N_2O + 2H_2O}\)
\end{enumerate}
\end{minipage}\\
\hline
\end{tabular}

\vspace{0.2cm}
\noindent \begin{tabular}{|p{0.9\linewidth}|}
\hline
\vspace{0.2cm}
\textbf{(23)} Balance: \(\mathrm{P_4 + O_2 \rightarrow P_4O_{10}}\)\\
\begin{minipage}[t]{0.85\linewidth}
\begin{enumerate}[label=\alph*.]
   \item \(\mathrm{P_4 + 10O_2 \rightarrow P_4O_{10}}\)
   \item \(\mathrm{P_4 + \tfrac{5}{2}O_2 \rightarrow P_4O_{10}}\)
   \item \(\mathrm{2P_4 + 5O_2 \rightarrow 2P_4O_{10}}\)
   \item \(\mathrm{P_4 + 5O_2 \rightarrow P_4O_{10}}\)
\end{enumerate}
\end{minipage}\\
\hline
\end{tabular}

\vspace{0.2cm}
\noindent \begin{tabular}{|p{0.9\linewidth}|}
\hline
\vspace{0.2cm}
\textbf{(24)} Balance: \(\mathrm{KMnO_4 + HCl \rightarrow KCl + MnCl_2 + Cl_2 + H_2O}\) (acidic chloride reaction)\\
\begin{minipage}[t]{0.85\linewidth}
\begin{enumerate}[label=\alph*.]
   \item \(\mathrm{KMnO_4 + 16HCl \rightarrow KCl + MnCl_2 + 5Cl_2 + 8H_2O}\)
   \item \(\mathrm{2KMnO_4 + 8HCl \rightarrow 2KCl + 2MnCl_2 + 5Cl_2 + 8H_2O}\)
   \item \(\mathrm{KMnO_4 + 8HCl \rightarrow KCl + MnCl_2 + 4Cl_2 + 4H_2O}\)
   \item \(\mathrm{2KMnO_4 + 16HCl \rightarrow 2KCl + 2MnCl_2 + 5Cl_2 + 8H_2O}\)
\end{enumerate}
\end{minipage}\\
\hline
\end{tabular}
\newpage

\vspace{0.2cm}
\noindent \begin{tabular}{|p{0.9\linewidth}|}
\hline
\vspace{0.2cm}
\textbf{(25)} Balance: \(\mathrm{C_6H_6 + O_2 \rightarrow CO_2 + H_2O}\) (benzene combustion)\\
\begin{minipage}[t]{0.85\linewidth}
\begin{enumerate}[label=\alph*.]
   \item \(\mathrm{C_6H_6 + 9O_2 \rightarrow 6CO_2 + 6H_2O}\)
   \item \(\mathrm{C_6H_6 + 3O_2 \rightarrow 6CO_2 + 3H_2O}\)
   \item \(\mathrm{2C_6H_6 + 15O_2 \rightarrow 12CO_2 + 6H_2O}\)
   \item \(\mathrm{C_6H_6 + 7.5O_2 \rightarrow 6CO_2 + 3H_2O}\)
\end{enumerate}
\end{minipage}\\
\hline
\end{tabular}

\vspace{0.2cm}
\noindent \begin{tabular}{|p{0.9\linewidth}|}
\hline
\vspace{0.2cm}
\textbf{(26)} Balance: \(\mathrm{H_2SO_4 + NaOH \rightarrow Na_2SO_4 + H_2O}\)\\
\begin{minipage}[t]{0.85\linewidth}
\begin{enumerate}[label=\alph*.]
   \item \(\mathrm{2H_2SO_4 + 2NaOH \rightarrow Na_2SO_4 + 2H_2O}\)
   \item \(\mathrm{H_2SO_4 + NaOH \rightarrow Na_2SO_4 + H_2O}\)
   \item \(\mathrm{H_2SO_4 + 2NaOH \rightarrow 2Na_2SO_4 + H_2O}\)
   \item \(\mathrm{H_2SO_4 + 2NaOH \rightarrow Na_2SO_4 + 2H_2O}\)
\end{enumerate}
\end{minipage}\\
\hline
\end{tabular}

\vspace{0.2cm}
\noindent \begin{tabular}{|p{0.9\linewidth}|}
\hline
\vspace{0.2cm}
\textbf{(27)} Balance: \(\mathrm{Al + HCl \rightarrow AlCl_3 + H_2}\)\\
\begin{minipage}[t]{0.85\linewidth}
\begin{enumerate}[label=\alph*.]
   \item \(\mathrm{Al + HCl \rightarrow AlCl_3 + H_2}\)
   \item \(\mathrm{Al + 3HCl \rightarrow AlCl_3 + \tfrac{3}{2}H_2}\)
   \item \(\mathrm{4Al + 12HCl \rightarrow 4AlCl_3 + 6H_2}\)
   \item \(\mathrm{2Al + 6HCl \rightarrow 2AlCl_3 + 3H_2}\)
\end{enumerate}
\end{minipage}\\
\hline
\end{tabular}

\vspace{0.2cm}
\noindent \begin{tabular}{|p{0.9\linewidth}|}
\hline
\vspace{0.2cm}
\textbf{(28)} Balance: \(\mathrm{FeS_2 + O_2 \rightarrow Fe_2O_3 + SO_2}\) (pyrite oxidation)\\
\begin{minipage}[t]{0.85\linewidth}
\begin{enumerate}[label=\alph*.]
   \item \(\mathrm{FeS_2 + O_2 \rightarrow Fe_2O_3 + SO_2}\)
   \item \(\mathrm{4FeS_2 + 5O_2 \rightarrow 2Fe_2O_3 + 8SO_2}\)
   \item \(\mathrm{4FeS_2 + 11O_2 \rightarrow 2Fe_2O_3 + 8SO_2}\)
   \item \(\mathrm{2FeS_2 + 5.5O_2 \rightarrow Fe_2O_3 + 4SO_2}\)
\end{enumerate}
\end{minipage}\\
\hline
\end{tabular}
\newpage

\vspace{0.2cm}
\noindent \begin{tabular}{|p{0.9\linewidth}|}
\hline
\vspace{0.2cm}
\textbf{(29)} Balance: \(\mathrm{C_4H_{10} + O_2 \rightarrow CO_2 + H_2O}\) (butane combustion)\\
\begin{minipage}[t]{0.85\linewidth}
\begin{enumerate}[label=\alph*.]
   \item \(\mathrm{C_4H_{10} + 9O_2 \rightarrow 4CO_2 + 5H_2O}\)
   \item \(\mathrm{C_4H_{10} + \tfrac{13}{2}O_2 \rightarrow 4CO_2 + 5H_2O}\)
   \item \(\mathrm{C_4H_{10} + 6.5O_2 \rightarrow 4CO_2 + 5H_2O}\)
   \item \(\mathrm{2C_4H_{10} + 13O_2 \rightarrow 8CO_2 + 10H_2O}\)
\end{enumerate}
\end{minipage}\\
\hline
\end{tabular}

\vspace{0.2cm}
\noindent \begin{tabular}{|p{0.9\linewidth}|}
\hline
\vspace{0.2cm}
\textbf{(30)} Balance: \(\mathrm{C_4H_8 + Br_2 \rightarrow C_4H_8Br_2}\) (addition)\\
\begin{minipage}[t]{0.85\linewidth}
\begin{enumerate}[label=\alph*.]
   \item \(\mathrm{C_4H_8 + 2Br_2 \rightarrow C_4H_8Br_4}\)
   \item \(\mathrm{C_4H_8 + Br_2 \rightarrow C_4H_8Br_2}\)
   \item \(\mathrm{2C_4H_8 + Br_2 \rightarrow 2C_4H_8Br_2}\)
   \item \(\mathrm{C_4H_8 + \tfrac{1}{2}Br_2 \rightarrow C_4H_8Br}\)
\end{enumerate}
\end{minipage}\\
\hline
\end{tabular}

\vspace{0.2cm}
\noindent \begin{tabular}{|p{0.9\linewidth}|}
\hline
\vspace{0.2cm}
\textbf{(31)} Balance: $\mathrm{H_2 + O_2 \rightarrow H_2O}$\\
\begin{minipage}[t]{0.85\linewidth}
\begin{enumerate}[label=\alph*.]
\item $\mathrm{H_2 + O_2 \rightarrow H_2O}$
\item $\mathrm{2H_2 + O_2 \rightarrow 2H_2O}$
\item $\mathrm{H_2 + 2O_2 \rightarrow 2H_2O}$
\item $\mathrm{4H_2 + O_2 \rightarrow 2H_2O}$
\end{enumerate}
\end{minipage}\\
\hline
\end{tabular}

\vspace{0.2cm}
\noindent \begin{tabular}{|p{0.9\linewidth}|}
\hline
\vspace{0.2cm}
\textbf{(32)} True or False: The coefficients in a balanced equation represent the ratio of atoms, not moles.\\
\textbf{Answer:} \underline{\hspace{2cm}}\\
\hline
\end{tabular}
\newpage

\vspace{0.2cm}
\noindent \begin{tabular}{|p{0.9\linewidth}|}
\hline
\vspace{0.2cm}
\textbf{(33)} How many moles of $\mathrm{CO_2}$ are in 88 g of $\mathrm{CO_2}$?\\
\begin{minipage}[t]{0.85\linewidth}
\begin{enumerate}[label=\alph*.]
\item 1 mol
\item 2 mol
\item 3 mol
\item 4 mol
\end{enumerate}
\end{minipage}\\
\hline
\end{tabular}

\vspace{0.2cm}
\noindent \begin{tabular}{|p{0.9\linewidth}|}
\hline
\vspace{0.2cm}
\textbf{(34)} True or False: The total mass of products must equal the total mass of reactants in a chemical equation.\\
\textbf{Answer:} \underline{\hspace{2cm}}\\
\hline
\end{tabular}

\vspace{0.2cm}
\noindent \begin{tabular}{|p{0.9\linewidth}|}
\hline
\vspace{0.2cm}
\textbf{(35)} Balance: $\mathrm{Al + O_2 \rightarrow Al_2O_3}$\\
\begin{minipage}[t]{0.85\linewidth}
\begin{enumerate}[label=\alph*.]
\item $\mathrm{2Al + 3O_2 \rightarrow Al_2O_3}$
\item $\mathrm{4Al + 3O_2 \rightarrow 2Al_2O_3}$
\item $\mathrm{Al + O_2 \rightarrow Al_2O_3}$
\item $\mathrm{3Al + 2O_2 \rightarrow Al_2O_3}$
\end{enumerate}
\end{minipage}\\
\hline
\end{tabular}

\vspace{0.2cm}
\noindent \begin{tabular}{|p{0.9\linewidth}|}
\hline
\vspace{0.2cm}
\textbf{(36)} True or False: The molar mass of a substance links the number of moles to the mass in grams.\\
\textbf{Answer:} \underline{\hspace{2cm}}\\
\hline
\end{tabular}
\newpage

\vspace{0.2cm}
\noindent \begin{tabular}{|p{0.9\linewidth}|}
\hline
\vspace{0.2cm}
\textbf{(37)} What is the molar mass of $\mathrm{CaCl_2}$?\\
\begin{minipage}[t]{0.85\linewidth}
\begin{enumerate}[label=\alph*.]
\item 75.5 g/mol
\item 95.0 g/mol
\item 110.98 g/mol
\item 147.0 g/mol
\end{enumerate}
\end{minipage}\\
\hline
\end{tabular}

\vspace{0.2cm}
\noindent \begin{tabular}{|p{0.9\linewidth}|}
\hline
\vspace{0.2cm}
\textbf{(38)} True or False: Moles and atoms are directly proportional.\\
\textbf{Answer:} \underline{\hspace{2cm}}\\
\hline
\end{tabular}

\vspace{0.2cm}
\noindent \begin{tabular}{|p{0.9\linewidth}|}
\hline
\vspace{0.2cm}
\textbf{(39)} How many atoms are in 2.0 mol of Na?\\
\begin{minipage}[t]{0.85\linewidth}
\begin{enumerate}[label=\alph*.]
\item $3.01\times10^{23}$
\item $6.02\times10^{23}$
\item $1.20\times10^{24}$
\item $2.40\times10^{24}$
\end{enumerate}
\end{minipage}\\
\hline
\end{tabular}

\vspace{0.2cm}
\noindent \begin{tabular}{|p{0.9\linewidth}|}
\hline
\vspace{0.2cm}
\textbf{(40)} True or False: A formula unit applies to ionic compounds.\\
\textbf{Answer:} \underline{\hspace{2cm}}\\
\hline
\end{tabular}
\newpage

\vspace{0.2cm}
\noindent \begin{tabular}{|p{0.9\linewidth}|}
\hline
\vspace{0.2cm}
\textbf{(41)} Convert 0.5 mol of $\mathrm{H_2O}$ to molecules.\\
\begin{minipage}[t]{0.85\linewidth}
\begin{enumerate}[label=\alph*.]
\item $1.20\times10^{23}$
\item $3.01\times10^{23}$
\item $6.02\times10^{23}$
\item $1.20\times10^{24}$
\end{enumerate}
\end{minipage}\\
\hline
\end{tabular}

\vspace{0.2cm}
\noindent \begin{tabular}{|p{0.9\linewidth}|}
\hline
\vspace{0.2cm}
\textbf{(42)} Balance: $\mathrm{N_2 + H_2 \rightarrow NH_3}$\\
\begin{minipage}[t]{0.85\linewidth}
\begin{enumerate}[label=\alph*.]
\item $\mathrm{N_2 + 3H_2 \rightarrow 2NH_3}$
\item $\mathrm{N_2 + H_2 \rightarrow NH_3}$
\item $\mathrm{2N_2 + 3H_2 \rightarrow 2NH_3}$
\item $\mathrm{3N_2 + 2H_2 \rightarrow NH_3}$
\end{enumerate}
\end{minipage}\\
\hline
\end{tabular}

\vspace{0.2cm}
\noindent \begin{tabular}{|p{0.9\linewidth}|}
\hline
\vspace{0.2cm}
\textbf{(43)} True or False: Subscripts can be changed to balance an equation.\\
\textbf{Answer:} \underline{\hspace{2cm}}\\
\hline
\end{tabular}

\vspace{0.2cm}
\noindent \begin{tabular}{|p{0.9\linewidth}|}
\hline
\vspace{0.2cm}
\textbf{(44)} Find the formula mass of $\mathrm{H_2SO_4}$.\\
\begin{minipage}[t]{0.85\linewidth}
\begin{enumerate}[label=\alph*.]
\item 49.0 g/mol
\item 75.0 g/mol
\item 98.1 g/mol
\item 120.0 g/mol
\end{enumerate}
\end{minipage}\\
\hline
\end{tabular}
\newpage

\vspace{0.2cm}
\noindent \begin{tabular}{|p{0.9\linewidth}|}
\hline
\vspace{0.2cm}
\textbf{(45)} True or False: Moles can be converted to grams using molar mass.\\
\textbf{Answer:} \underline{\hspace{2cm}}\\
\hline
\end{tabular}

\vspace{0.2cm}
\noindent \begin{tabular}{|p{0.9\linewidth}|}
\hline
\vspace{0.2cm}
\textbf{(46)} Balance: $\mathrm{CH_4 + O_2 \rightarrow CO_2 + H_2O}$\\
\begin{minipage}[t]{0.85\linewidth}
\begin{enumerate}[label=\alph*.]
\item $\mathrm{CH_4 + 2O_2 \rightarrow CO_2 + 2H_2O}$
\item $\mathrm{2CH_4 + O_2 \rightarrow 2CO_2 + H_2O}$
\item $\mathrm{CH_4 + O_2 \rightarrow CO + H_2O}$
\item $\mathrm{2CH_4 + 4O_2 \rightarrow 4CO_2 + 8H_2O}$
\end{enumerate}
\end{minipage}\\
\hline
\end{tabular}

\vspace{0.2cm}
\noindent \begin{tabular}{|p{0.9\linewidth}|}
\hline
\vspace{0.2cm}
\textbf{(47)} True or False: The coefficient ``2'' in $2\mathrm{H_2O}$ means there are 4 hydrogen atoms and 2 oxygen atoms.\\
\textbf{Answer:} \underline{\hspace{2cm}}\\
\hline
\end{tabular}

\vspace{0.2cm}
\noindent \begin{tabular}{|p{0.9\linewidth}|}
\hline
\vspace{0.2cm}
\textbf{(48)} How many moles of $\mathrm{NaCl}$ are in 117 g of $\mathrm{NaCl}$?\\
\begin{minipage}[t]{0.85\linewidth}
\begin{enumerate}[label=\alph*.]
\item 1.0 mol
\item 2.0 mol
\item 3.0 mol
\item 4.0 mol
\end{enumerate}
\end{minipage}\\
\hline
\end{tabular}
\newpage

\vspace{0.2cm}
\noindent \begin{tabular}{|p{0.9\linewidth}|}
\hline
\vspace{0.2cm}
\textbf{(49)} True or False: Formula mass is the sum of atomic masses in a compound.\\
\textbf{Answer:} \underline{\hspace{2cm}}\\
\hline
\end{tabular}

\vspace{0.2cm}
\noindent \begin{tabular}{|p{0.9\linewidth}|}
\hline
\vspace{0.2cm}
\textbf{(50)} Balance: $\mathrm{Fe + O_2 \rightarrow Fe_2O_3}$\\
\begin{minipage}[t]{0.85\linewidth}
\begin{enumerate}[label=\alph*.]
\item $\mathrm{4Fe + 3O_2 \rightarrow 2Fe_2O_3}$
\item $\mathrm{2Fe + O_2 \rightarrow Fe_2O_3}$
\item $\mathrm{3Fe + 2O_2 \rightarrow Fe_2O_3}$
\item $\mathrm{Fe + O_2 \rightarrow Fe_2O_3}$
\end{enumerate}
\end{minipage}\\
\hline
\end{tabular}

\vspace{0.2cm}
\noindent \begin{tabular}{|p{0.9\linewidth}|}
\hline
\vspace{0.2cm}
\textbf{(51)} True or False: $6.02\times10^{23}$ formula units equal one mole of an ionic compound.\\
\textbf{Answer:} \underline{\hspace{2cm}}\\
\hline
\end{tabular}

\vspace{0.2cm}
\noindent \begin{tabular}{|p{0.9\linewidth}|}
\hline
\vspace{0.2cm}
\textbf{(52)} Convert 3.0 mol of $\mathrm{O_2}$ to molecules.\\
\begin{minipage}[t]{0.85\linewidth}
\begin{enumerate}[label=\alph*.]
\item $1.81\times10^{23}$
\item $6.02\times10^{23}$
\item $1.81\times10^{24}$
\item $3.61\times10^{24}$
\end{enumerate}
\end{minipage}\\
\hline
\end{tabular}
\newpage

\vspace{0.2cm}
\noindent \begin{tabular}{|p{0.9\linewidth}|}
\hline
\vspace{0.2cm}
\textbf{(53)} True or False: One mole of $\mathrm{H_2O}$ contains the same number of oxygen atoms as one mole of $\mathrm{CO_2}$.\\
\textbf{Answer:} \underline{\hspace{2cm}}\\
\hline
\end{tabular}

\vspace{0.2cm}
\noindent \begin{tabular}{|p{0.9\linewidth}|}
\hline
\vspace{0.2cm}
\textbf{(54)} What is the molar mass of $\mathrm{Mg(OH)_2}$?\\
\begin{minipage}[t]{0.85\linewidth}
\begin{enumerate}[label=\alph*.]
\item 34.0 g/mol
\item 48.3 g/mol
\item 58.3 g/mol
\item 74.3 g/mol
\end{enumerate}
\end{minipage}\\
\hline
\end{tabular}

\vspace{0.2cm}
\noindent \begin{tabular}{|p{0.9\linewidth}|}
\hline
\vspace{0.2cm}
\textbf{(55)} True or False: Balancing equations ensures conservation of charge as well as mass.\\
\textbf{Answer:} \underline{\hspace{2cm}}\\
\hline
\end{tabular}

\vspace{0.2cm}
\noindent \begin{tabular}{|p{0.9\linewidth}|}
\hline
\vspace{0.2cm}
\textbf{(56)} Balance: $\mathrm{K + Cl_2 \rightarrow KCl}$\\
\begin{minipage}[t]{0.85\linewidth}
\begin{enumerate}[label=\alph*.]
\item $\mathrm{K + Cl_2 \rightarrow KCl}$
\item $\mathrm{2K + Cl_2 \rightarrow 2KCl}$
\item $\mathrm{2K + 2Cl_2 \rightarrow 2KCl_2}$
\item $\mathrm{K + 2Cl_2 \rightarrow 2KCl}$
\end{enumerate}
\end{minipage}\\
\hline
\end{tabular}
\newpage

\vspace{0.2cm}
\noindent \begin{tabular}{|p{0.9\linewidth}|}
\hline
\vspace{0.2cm}
\textbf{(57)} True or False: The number of atoms of each element must be equal on both sides of a balanced equation.\\
\textbf{Answer:} \underline{\hspace{2cm}}\\
\hline
\end{tabular}

\vspace{0.2cm}
\noindent \begin{tabular}{|p{0.9\linewidth}|}
\hline
\vspace{0.2cm}
\textbf{(58)} How many grams are in 0.25 mol of $\mathrm{Na_2CO_3}$?\\
\begin{minipage}[t]{0.85\linewidth}
\begin{enumerate}[label=\alph*.]
\item 13.25 g
\item 26.5 g
\item 26.8 g
\item 53.0 g
\end{enumerate}
\end{minipage}\\
\hline
\end{tabular}

\vspace{0.2cm}
\noindent \begin{tabular}{|p{0.9\linewidth}|}
\hline
\vspace{0.2cm}
\textbf{(59)} True or False: One mole of hydrogen gas and one mole of oxygen gas contain the same number of molecules.\\
\textbf{Answer:} \underline{\hspace{2cm}}\\
\hline
\end{tabular}

\vspace{0.2cm}
\noindent \begin{tabular}{|p{0.9\linewidth}|}
\hline
\vspace{0.2cm}
\textbf{(60)} Balance: $\mathrm{Na + H_2O \rightarrow NaOH + H_2}$\\
\begin{minipage}[t]{0.85\linewidth}
\begin{enumerate}[label=\alph*.]
\item $\mathrm{Na + H_2O \rightarrow NaOH + H_2}$
\item $\mathrm{2Na + 2H_2O \rightarrow 2NaOH + H_2}$
\item $\mathrm{Na + 2H_2O \rightarrow NaOH + 2H_2}$
\item $\mathrm{2Na + H_2O \rightarrow NaOH + 2H_2}$
\end{enumerate}
\end{minipage}\\
\hline
\end{tabular}

\newpage
{\large\textbf{Checking for Understanding (general conceptual questions)}}\\
\noindent \begin{tabular}{|p{0.9\linewidth}|}
\hline
\vspace{0.2cm}

\begin{enumerate}
    \item \textbf{Targeted focus through self-selection:} When students choose which questions to tackle from the worksheet, they naturally gravitate toward areas they feel less confident about, ensuring the group addresses gaps in understanding.
    \item \textbf{Peer modeling of problem-solving:} Seeing a peer work through a problem on the board exposes students to different ways of setting up and explaining a solution, clarifying concepts for observers.
    \item \textbf{Reflection on understanding:} Students must decide whether they can solve a question themselves or need help, prompting metacognition and highlighting what they truly understand versus what requires clarification.
    \item \textbf{Immediate correction of misconceptions:} Peers correcting or expanding an explanation in real time prevents misunderstandings from becoming ingrained.
    \item \textbf{Balancing independence and collaboration:} Students first work independently on a question, then collaborate during discussion, ensuring individual accountability while leveraging group problem-solving.
    \item \textbf{Distinguishing conceptual vs. procedural questions:} Students can classify questions by identifying key ideas versus stepwise calculations, determining whether a question emphasizes big-picture reasoning or step-by-step execution.
    \item \textbf{Connecting big concepts to details:} Rotating through different student-selected questions helps the group link overarching concepts to multiple specific examples, reinforcing understanding of the framework behind the problems.
    \item \textbf{Verbalizing reasoning:} Explaining steps out loud forces students to articulate assumptions, logic, and chemical principles, improving clarity and communication of ideas.
    \item \textbf{Learning from multiple approaches:} Different students may solve the same problem uniquely, exposing peers to alternative methods or perspectives and broadening their problem-solving toolkit.
    \item \textbf{Application of theory to practice:} Working collaboratively on worksheet problems allows students to connect abstract chemical principles to concrete calculations or predictions.
    \item \textbf{Instructor insight:} Questions students select or ask reveal where common difficulties lie, helping instructors or SI leaders provide targeted guidance.
    \item \textbf{Strengthening long-term understanding:} Actively solving, explaining, and correcting questions in real time consolidates memory and understanding more effectively than passive review.
\end{enumerate}
\hline
\end{tabular}

\newpage
\noindent
\newpage
\noindent
\section*{\underline{Activity 2}}
\begin{tabular}{|p{4.2cm}|p{9.8cm}|}
\hline
\textbf{Collaborative Learning Techniques} & Clusters \\
\hline
\textbf{Learning Strategy} & 3 Before Me \\
\hline
\textbf{Time} & 12:35---12:50 \\
\hline
\textbf{Materials} & A computer, OWLv2, textbook, Canvas, Anki, Anki decks for download, readiness to speak or share ideas \\
\hline
\textbf{Lecture/Textbook/Old Plans/Slide References} & Anki, Canvas/Cengage, Chapters 1--6 \\
\hline
\end{tabular}

\vspace{0.5cm}
\noindent
\begin{tabular}{|p{0.9\linewidth}|}
\hline
{\large\textbf{Definitions}}
\begin{enumerate}
    \item \textbf{3 Before Me:} When a student asks a question during a session, have 3 students (or less depending on the number of students in the session) comment on a unique feature of that idea. The SI leader will mediate correct responses and help fill in gaps in understanding. This is a good strategy to model at the beginning of the semester and use throughout. This can help with redirecting questions and to encourage student-to-student interactions.
    \end{itemize}
\end{enumerate} \\
\hline
\end{tabular}
\newpage

\vspace{0.5cm}
\noindent
{\large \textbf{Instructions (detailed activity breakdown below)}} \\

\begin{tabular}{|p{0.9\linewidth}|}
\hline
\begin{enumerate}
    \item \textbf{Purpose:} The goal of the 3 Before Me strategy is to encourage students to collaborate, listen actively to peers, and learn to evaluate and build upon others’ reasoning before relying on the instructor or SI leader for answers.
    
    \item \textbf{Setup (2 minutes):} Arrange students into clusters of 3--5. Make sure all students can see shared materials (slides, textbook examples, or OWLv2 problems). Explain that before any question is answered by the SI leader, at least three students must contribute their thoughts, ideas, or possible solutions.
    
    \item \textbf{Introducing the activity (2 minutes):} The SI leader explains the rule clearly: when someone asks a question, three peers must respond first. Responses can include restating the problem, identifying known information, describing what concept applies, or explaining a partial solution.
    
    \item \textbf{Running the activity (8 minutes):}
    \begin{itemize}
        \item A student poses a question related to the lecture, homework, or Anki review material.
        \item The SI leader redirects the question to the group, saying, “Let’s hear three ideas before I step in.”
        \item Three (or fewer, depending on group size) students share distinct perspectives, explanations, or steps toward a solution.
        \item The SI leader mediates, affirming correct reasoning and gently addressing misconceptions by asking probing questions like “Can someone add to that?” or “Why might that be true?”
        \item If needed, the SI leader provides clarification or synthesis after the three contributions.
    \end{itemize}
    
    \item \textbf{Role of the SI Leader:} 
    \begin{itemize}
        \item Model active listening and respectful engagement.
        \item Redirect questions that go directly to the leader back to peers (“What do others think?”).
        \item Encourage concise, evidence-based explanations.
        \item Highlight when peer responses align with lecture or textbook content.
    \end{itemize}
    
    \item \textbf{Variation:} If participation is slow, allow students to discuss in pairs for 30 seconds before sharing aloud. Alternatively, assign roles (Concept Clarifier, Example Provider, Connector) to prompt varied responses.
    
    \item \textbf{Wrap-up (3 minutes):}
    \begin{itemize}
        \item Summarize key takeaways from the questions discussed.
        \item Highlight moments when students corrected or built upon each other’s reasoning.
        \item Encourage reflection: ask, “What did you learn from hearing others’ explanations?”
        \item Remind students that this approach mirrors how scientists collaborate and refine understanding.
    \end{itemize}
    
    \item \textbf{Outcome:} By the end of the activity, students practice peer instruction, collaborative reasoning, and self-correction, while the SI leader gathers insight into where conceptual gaps remain.
\end{enumerate} \\
\hline
\end{tabular}

\newpage
\noindent
{\large\textbf{Checking for Understanding (questions \& answers)}}\\[0.2em]
\noindent
\begin{tabular}{|p{0.9\linewidth}|}
\hline
\textbf{Questions:} \\
\begin{enumerate}
    \item \textbf{Targeted clarification:} Allowing students to bring questions from the lecture lets them focus on concepts they found most confusing.
    \item \textbf{Peer modeling of explanations:} Observing how classmates summarize or interpret lecture content helps students see alternative ways to understand the material.
    \item \textbf{Reflection on comprehension:} Discussing lecture points collaboratively encourages students to distinguish what they grasp from what needs further clarification.
    \item \textbf{Immediate correction of misconceptions:} Peers can correct or expand on explanations in real time, preventing misunderstandings from solidifying.
    \item \textbf{Balancing independent thinking and discussion:} Students first reflect individually, then contribute to group discussion, promoting both accountability and collaborative learning.
    \item \textbf{Identifying conceptual vs. procedural points:} Students can highlight which parts of the lecture are big-picture ideas versus detailed steps or examples.
    \item \textbf{Connecting concepts across topics:} Discussing multiple lecture points together helps students link overarching principles to specific details.
    \item \textbf{Verbalizing reasoning:} Explaining ideas out loud strengthens students’ ability to communicate and reason about the material clearly.
    \item \textbf{Learning from diverse perspectives:} Hearing multiple interpretations or summaries of the same lecture point exposes students to alternative approaches.
    \item \textbf{Application of theory:} Collaborative discussion encourages students to connect lecture content to examples, problems, or real-world applications.
    \item \textbf{Instructor insight:} Observing which lecture points students question most helps the instructor identify common difficulties and adjust instruction.
    \item \textbf{Reinforcing long-term understanding:} Actively reviewing and discussing lecture material together solidifies comprehension better than passive note review.
\end{enumerate} \\
\hline  
\end{tabular}

\newpage
\noindent
\section*{\underline{Closer}}
\begin{tabular}{|p{4.2cm}|p{9.8cm}|}
\hline
\textbf{Collaborative Learning Techniques} & Group Survey\\
\hline
\textbf{Learning Strategy} & Brain dump\\
\hline
\textbf{Time} & 12:50---12:55\\
\hline
\textbf{Materials} & Index cards, writing utensils\\
\hline
\textbf{Lecture/Textbook/Old Plans/Slide References} & Chapter 3: Periodic Table, Chapter 4: Chemical Bonding: The Ionic Bond Model, Chapter 5: The Chemical Bonding: The Covalent Bond Model \\
\hline
\end{tabular}

\vspace{0.5cm}
\noindent
\begin{tabular}{|p{0.9\linewidth}|}
\hline
{\large\textbf{Definitions}}
\begin{enumerate}
   \item\textbf{Group survey:} Each group member is surveyed to discover their position on an issue, problem or topic. This process ensures that each member of the group is allowed to offer or state their point of view.
   \item\textbf{Brain dump:} This is a great closer for students to sum up their knowledge on the topic(s) from the session. Have the students take out a blank piece of paper /Google Doc and write/type("dump") everything they know relevant from the discussion that day. The SI Leader can look at their work to assess the students' understanding of the course material and have them compare it to each other. It is important to at least discuss with the students so the SI Leader knows the students' comfort levels with the material. 
\end{enumerate}
\\ \hline
\end{tabular}

\newpage
\noindent
{\large\textbf{Instructions (please provide a\hl{DETAILED} activity breakdown below (and/or attached}}

\begin{tabular}{|p{0.9\linewidth}|}
\hline
\begin{enumerate}[itemsep=0.3em]
   \item Students are put into small groups and are asked to write down their confidence on a single notecard without names.
   \item The cards are evidence of the group's confidence in the content.
\end{enumerate}
\\ \hline
\end{tabular}

\vspace{0.5cm}
\noindent
{\large\textbf{Content (fully worked out problems \& answers)}

\begin{tabular}{|p{0.9\linewidth}|}
\hline
\begin{itemize}
   \item No new worked problems; brain dump.
\end{itemize}
\\ \hline
\end{tabular}

\vspace{0.5cm}
\noindent
{\large\textbf{Checking for Understanding (questions \& answers) continued...}

\begin{tabular}{|p{0.9\linewidth}|}
\hline
\begin{itemize}
   \item Review confidence levels
   \item Interpret the variety of confidence and their distribution in groups.
   \item Get feedback for what they like and don't like
   \item Check in if the time is suitable for the students
\end{itemize}
\\ \hline
\end{tabular}
\newpage
\section*{Plan Checklist}

\begin{tabular}{|p{0.9\linewidth}|}
\hline
\begin{itemize}
\item[$\Box$] Variety of Collaborative Learning Techniques and Learning Strategies
\item[$\Box$] Chapter/slide\# where information is coming from
\item[$\Box$] Included timestamps (and buffer time as needed)
\item[$\Box$] Details on how leader will break up students
\item[$\Box$] Checking for Understanding questions
\item[$\Box$] Fully worked out problems and examples for each activity
\item[$\Box$] Necessary links required for online implementation (if applicable)
\item[$\Box$] Source material correctly cited (if applicable)
\item[$\Box$] Plan is uploaded to Teams
\end{itemize}
\\ \hline
\end{tabular}
\newpage

\end{document}
